% GENERATED FILE. Edit canonical lessons, then run npm run generate:books.
% canonical-source-hash: fnv1a64:a69095f70d7646f7
% canonical-lessons: KA-C274-kuni, KA-C274-ottu, KA-C274-jama-madu, KA-C274-himpade, KA-C274-sahi-madu

\chapter{To dance, To press, To deposit, To withdraw, To sign}
\label{ch:ka-to-dance-to-press-to-deposit-to-withdraw-to-sign}
\hlchaptermodality{274}

\begin{chapteropening}
\textbf{By the end of this chapter:} \emph{I can say to dance, to press, to deposit, to withdraw and to sign.}

Complete the last lesson of chapter 274: I can say to dance, to press, to deposit, to withdraw and to sign.
\end{chapteropening}

\section[\texorpdfstring{\kn{ಕುಣಿ} (kuṇi)}{ಕುಣಿ (kuṇi)}]{\kn{ಕುಣಿ} (kuṇi) — to dance}

\label{lesson:KA-C274-kuni}



\begin{quote}
Before the new one: say the Kannada for to drive, then the Kannada for to slip.
\end{quote}

\subsection*{You'll want to know: \kn{ಕುಣಿ}}
\textbf{\kn{ಕುಣಿ}} — \emph{kuṇi} — \textquotedblleft{}to dance\textquotedblright{}.

\begin{grammarlens}[title={What you've built}]
One more word for this chapter.
\end{grammarlens}

\subsection*{Your turn}
\begin{itemize}
\raggedright
  \item \emph{Say it:} \emph{kuṇi}
  \item \emph{Say it:} \emph{kuṇi}, once more
  \item \emph{Say it:} say \emph{kuṇi}
  \item \emph{Recall:} say the Kannada for to drive, then the Kannada for to slip, then say \emph{kuṇi} again
\end{itemize}

\begin{tcolorbox}[breakable,colback=teal!4,colframe=teal!35!black,title={Before you move on}]
What does \kn{ಕುಣಿ} mean? (\textquotedblleft{}To dance\textquotedblright{}.) Say it once more.
\end{tcolorbox}
\section[\texorpdfstring{\kn{ಒತ್ತು} (ottu)}{ಒತ್ತು (ottu)}]{\kn{ಒತ್ತು} (ottu) — to press}

\label{lesson:KA-C274-ottu}



\begin{quote}
Before the new one: say the Kannada for to slip, then the Kannada for to dance.
\end{quote}

\subsection*{You'll want to know: \kn{ಒತ್ತು}}
\textbf{\kn{ಒತ್ತು}} — \emph{ottu} — \textquotedblleft{}to press\textquotedblright{}.

\begin{grammarlens}[title={What you've built}]
One more word for this chapter.
\end{grammarlens}

\subsection*{Your turn}
\begin{itemize}
\raggedright
  \item \emph{Say it:} \emph{ottu}
  \item \emph{Say it:} \emph{ottu}, once more
  \item \emph{Say it:} say \emph{ottu}
  \item \emph{Recall:} say the Kannada for to slip, then the Kannada for to dance, then say \emph{ottu} again
\end{itemize}

\begin{tcolorbox}[breakable,colback=teal!4,colframe=teal!35!black,title={Before you move on}]
What does \kn{ಒತ್ತು} mean? (\textquotedblleft{}To press\textquotedblright{}.) Say it once more.
\end{tcolorbox}
\section[\texorpdfstring{\kn{ಜಮಾ} \kn{ಮಾಡು} (jamā māḍu)}{ಜಮಾ ಮಾಡು (jamā māḍu)}]{\kn{ಜಮಾ} \kn{ಮಾಡು} (jamā māḍu) — to deposit (money)}

\label{lesson:KA-C274-jama-madu}



\begin{quote}
Before the new one: say the Kannada for to dance, then the Kannada for to press.
\end{quote}

\subsection*{You'll want to know: \kn{ಜಮಾ} \kn{ಮಾಡು}}
\textbf{\kn{ಜಮಾ} \kn{ಮಾಡು}} — \emph{jamā māḍu} — \textquotedblleft{}to deposit (money)\textquotedblright{}.

\begin{grammarlens}[title={What you've built}]
One more word for this chapter.
\end{grammarlens}

\subsection*{Your turn}
\begin{itemize}
\raggedright
  \item \emph{Say it:} \emph{jamā māḍu}
  \item \emph{Say it:} \emph{jamā māḍu}, once more
  \item \emph{Say it:} say \emph{jamā māḍu}
  \item \emph{Recall:} say the Kannada for to dance, then the Kannada for to press, then say \emph{jamā māḍu} again
\end{itemize}

\begin{tcolorbox}[breakable,colback=teal!4,colframe=teal!35!black,title={Before you move on}]
What does \kn{ಜಮಾ} \kn{ಮಾಡು} mean? (\textquotedblleft{}To deposit (money)\textquotedblright{}.) Say it once more.
\end{tcolorbox}
\section[\texorpdfstring{\kn{ಹಿಂಪಡೆ} (himpaḍe)}{ಹಿಂಪಡೆ (himpaḍe)}]{\kn{ಹಿಂಪಡೆ} (himpaḍe) — to withdraw (money)}

\label{lesson:KA-C274-himpade}



\begin{quote}
Before the new one: say the Kannada for to press, then the Kannada for to deposit.
\end{quote}

\subsection*{You'll want to know: \kn{ಹಿಂಪಡೆ}}
\textbf{\kn{ಹಿಂಪಡೆ}} — \emph{himpaḍe} — \textquotedblleft{}to withdraw (money)\textquotedblright{}.

\begin{grammarlens}[title={What you've built}]
One more word for this chapter.
\end{grammarlens}

\subsection*{Your turn}
\begin{itemize}
\raggedright
  \item \emph{Say it:} \emph{himpaḍe}
  \item \emph{Say it:} \emph{himpaḍe}, once more
  \item \emph{Say it:} say \emph{himpaḍe}
  \item \emph{Recall:} say the Kannada for to press, then the Kannada for to deposit, then say \emph{himpaḍe} again
\end{itemize}

\begin{tcolorbox}[breakable,colback=teal!4,colframe=teal!35!black,title={Before you move on}]
What does \kn{ಹಿಂಪಡೆ} mean? (\textquotedblleft{}To withdraw (money)\textquotedblright{}.) Say it once more.
\end{tcolorbox}
\section[\texorpdfstring{\kn{ಸಹಿ} \kn{ಮಾಡು} (sahi māḍu)}{ಸಹಿ ಮಾಡು (sahi māḍu)}]{\kn{ಸಹಿ} \kn{ಮಾಡು} (sahi māḍu) — to sign}

\label{lesson:KA-C274-sahi-madu}



\begin{quote}
Before the new one: say the Kannada for to deposit, then the Kannada for to withdraw.
\end{quote}

\subsection*{You'll want to know: \kn{ಸಹಿ} \kn{ಮಾಡು}}
\textbf{\kn{ಸಹಿ} \kn{ಮಾಡು}} — \emph{sahi māḍu} — \textquotedblleft{}to sign\textquotedblright{}.

\begin{grammarlens}[title={What you've built}]
One more word for this chapter.
\end{grammarlens}

\subsection*{Your turn}
\begin{itemize}
\raggedright
  \item \emph{Say it:} \emph{sahi māḍu}
  \item \emph{Say it:} \emph{sahi māḍu}, once more
  \item \emph{Say it:} say \emph{sahi māḍu}
  \item \emph{Recall:} say the Kannada for to deposit, then the Kannada for to withdraw, then say \emph{sahi māḍu} again
\end{itemize}

\begin{tcolorbox}[breakable,colback=teal!4,colframe=teal!35!black,title={Before you move on}]
What does \kn{ಸಹಿ} \kn{ಮಾಡು} mean? (\textquotedblleft{}To sign\textquotedblright{}.) Say it once more.
\end{tcolorbox}
